\begin{corollary}[5.9.5]
\label{IV.2.s5.9.5-fr}
Bajo las hipótesis de \textup{(5.9.4)}, si
$\mathcal H^0_{X/Z}(\mathcal F)$ es coherente, lo mismo ocurre con
$\mathcal H^0_{X'/Z'}(f^*(\mathcal F))$. El recíproco es cierto cuando $f$
es un morfismo fielmente plano y cuasicompacto.
\end{corollary}

La primera afirmación resulta de (5.9.4.1) y de
$(\mathbf{0}_{\mathrm I},5.3.11)$; la segunda equivale a decir que, si
$\mathcal H^0_{X'/Z'}(f^*(\mathcal F))$ es de tipo finito, lo mismo ocurre
con $\mathcal H^0_{X/Z}(\mathcal F)$; esto resulta de (5.9.4.1) y de (2.5.2).

\begin{corollary}[5.9.6]
\label{IV.2.s5.9.6-fr}
Sean $X$ un preesquema localmente noetheriano, $Z$ una parte de $X$ estable
por especialización y $\mathcal F$ un $\mathcal O_X$-módulo cuasicoherente.
Para todo $x\in X$, pongamos $X_x=\operatorname{Spec}(\mathcal O_x)$,
$Z_x=Z\cap X_x$; se tiene un isomorfismo canónico funtorial
\begin{equation}
(\mathcal H^0_{X/Z}(\mathcal F))_x
\xrightarrow{\sim}\mathcal H^0_{X_x/Z_x}(\widetilde{\mathcal F}_x)
\tag{5.9.6.1}\label{IV.2.s5.9.6.1-fr}
\end{equation}
\end{corollary}

Basta aplicar (5.9.4) al morfismo canónico $X_x\to X$, que es plano, y tener
en cuenta (\textbf I, 1.6.5).

\begin{env}[5.9.7]
\label{IV.2.s5.9.7-fr}
Con las notaciones de (5.9.1), se tiene, para todo $\alpha$, un homomorfismo
canónico funtorial
$\mathcal F\longrightarrow(i_\alpha)_*(\mathcal F|U_\alpha)$
$(\mathbf{0}_{\mathrm I},4.4.3.2)$,
y estos homomorfismos forman un sistema inductivo; pasando al límite
inductivo, se deduce, pues, un homomorfismo canónico funtorial
\begin{equation}
\rho_{X/Z}:\mathcal F\longrightarrow\mathcal H^0_{X/Z}(\mathcal F)
\tag{5.9.7.1}\label{IV.2.s5.9.7.1-fr}
\end{equation}
\end{env}

\begin{proposition}[5.9.8]
\EGAReaderTarget{ega4.c04.prop.5-9-8}{IV.2.s5.9.8-fr}
Sean $X$ un preesquema localmente noetheriano, $Z$ una parte de $X$ estable
por especialización y $\mathcal F$ un $\mathcal O_X$-módulo. Las propiedades
siguientes son equivalentes:
\begin{enumerate}
\item[\textup{a)}] El homomorfismo $\rho_{X/Z}$ \textup{(5.9.7.1)} es inyectivo
(\textup{resp.} biyectivo).
\item[\textup{b)}] Para todo abierto noetheriano $V$ de $X$, el homomorfismo
\[
(\rho_{X/Z})_V:\Gamma(V,\mathcal F)\longrightarrow
\varinjlim_\alpha\Gamma(V\cap U_\alpha,\mathcal F)
\]
es inyectivo (\textup{resp.} biyectivo).
\item[\textup{a')}] Para toda parte cerrada $T\subset Z$ de $X$, el homomorfismo
canónico $(\mathbf{0}_{\mathrm I},4.4.3.2)$
\[
\mathcal F\longrightarrow i_*(\mathcal F|X\setmin T)
\]
(donde $i:X\setmin T\to X$ es la inyección canónica) es inyectivo
(\textup{resp.} biyectivo).
\item[\textup{b')}] Para toda parte cerrada $T\subset Z$ de $X$ y todo abierto
noetheriano $V$ de $X$, el homomorfismo de restricción
\[
\Gamma(V,\mathcal F)\longrightarrow
\Gamma(V\cap(X\setmin T),\mathcal F)
\]
es inyectivo (\textup{resp.} biyectivo).
\end{enumerate}
\end{proposition}

Teniendo en cuenta (5.9.1.2), la equivalencia entre \emph{a)} y \emph{b)}
(resp. entre \emph{a')} y \emph{b')}) se sigue de
\oldpage[IV-2]{112}

la definición del funtor $\Gamma$ y del hecho de que es exacto por la
izquierda. Como el homomorfismo $(\rho_{X/Z})_V$ es el compuesto
\begin{equation}
\Gamma(V,\mathcal F)\longrightarrow
\Gamma(V\cap(X\setmin T),\mathcal F)\longrightarrow
\varinjlim_\alpha\Gamma(V\cap U_\alpha,\mathcal F)
\tag{5.9.8.1}\label{IV.2.s5.9.8.1-fr}
\end{equation}
para toda parte cerrada $T\subset Z$, si $(\rho_{X/Z})_V$ es inyectivo, lo
mismo ocurre con
$\Gamma(V,\mathcal F)\to\Gamma(V\cap(X\setmin T),\mathcal F)$; por otra parte, el
hecho de que \emph{b')} implique \emph{b)} resulta de la definición de un límite inductivo.
Queda demostrar que, si $\rho_{X/Z}$ es biyectivo, lo mismo ocurre con
$\Gamma(V,\mathcal F)\to\Gamma(V\cap(X\setmin T),\mathcal F)$ y, para ello, basta,
en virtud de (5.9.8.1), ver que, si $U'\subset U$ son dos abiertos contenidos
en $V$ y que contienen $V\cap(X-Z)$, el homomorfismo de restricción
$\Gamma(U,\mathcal F)\to\Gamma(U',\mathcal F)$ es inyectivo; pero esto resulta
de que $\rho_{X/Z}$ es inyectivo, sustituyendo en lo que precede $V$ por $U$
y $V\cap(X\setmin T)$ por $U'$.
